\chapter{电磁环境怎样响应电子}
\label{chap:feqo-electromagnetic-field}

上一章把入射电子写成可归一化的正能波包。即使给定了电子电流，
我们仍不知道它在纳米结构附近激起什么电场：边界和材料会改变每个频率的响应。
Maxwell 方程让这个问题成为一个带边界条件的线性方程；
Green 张量把同一结构对任意电流的响应保存下来。

\section{由 Maxwell 方程构造响应核}
\label{sec:feqo-green}

全书把实场写成正、负频部分之和，频域因子取
\(\ee^{-\ii\omega t}\)。先假设材料在线性响应范围内，且在所研究的尺度
可用局域相对介电张量 \(\boldsymbol\varepsilon(\mathbf r,\omega)\)
和相对磁导率 \(\mu_r(\mathbf r,\omega)\) 描述。
给定外加电流 \(\mathbf j(\mathbf r,\omega)\)，
Maxwell 的两个旋度方程变为
\[
\nabla\times\mathbf E=\ii\omega\mu_0\mu_r\mathbf H,\qquad
\nabla\times\mathbf H
=\mathbf j-\ii\omega\varepsilon_0
\boldsymbol\varepsilon\mathbf E.
\]
第一式给
\(\mathbf H=(\ii\omega\mu_0)^{-1}\mu_r^{-1}\nabla\times\mathbf E\)。
代入第二式并用 \(c^{-2}=\varepsilon_0\mu_0\)，得到
\begin{equation}
\mathcal L_\omega\mathbf E
\equiv
\left[\nabla\times\mu_r^{-1}\nabla\times
-\frac{\omega^2}{c^2}\boldsymbol\varepsilon\right]\mathbf E
=\ii\omega\mu_0\mathbf j .
\label{eq:feqo-maxwell-operator}
\end{equation}
材料和几何进入左边的算子，电子轨迹进入右边的源。
我们希望为同一算子只求一次“单位点源的场”。

\begin{definition}[推迟 Green 张量]
\(\mathbf G(\mathbf r,\mathbf r';\omega)\) 是张量分布，满足
\begin{equation}
\mathcal L_\omega\mathbf G(\mathbf r,\mathbf r';\omega)
=\mathbf I\,\delta^{(3)}(\mathbf r-\mathbf r').
\label{eq:feqo-green-definition}
\end{equation}
它在无穷远取向外辐射条件，或在有限装置上取指定边界条件；
频率轴上的极点由 \(\omega\to\omega+\ii0\) 从上半平面逼近。
这些条件一同选出推迟解。
\end{definition}

“推迟”不是给核加的形容词。定义它的时间核为
\(\mathbf G(t)=\int\frac{\dd\omega}{2\pi}
\mathbf G(\omega)\ee^{-\ii\omega t}\)。
若 \(\mathbf G(\omega)\) 在上半复频平面解析，
且高频增长允许作相应收敛处理，那么 \(t<0\) 时可把
积分路径向上闭合：因 \(\ee^{-\ii\omega t}\) 在那里衰减，
闭合区域没有极点，积分为零。
因此先于电流源的时刻没有响应。
频率轴上的 \(+\ii0\) 正是在实轴上规定从哪一侧取边界值；
若误改成 \(-\ii0\)，时间支撑与辐射条件也会改变。

把式~\eqref{eq:feqo-green-definition} 乘源电流并对
\(\mathbf r'\) 积分，再利用 delta 分布，就有
\begin{equation}
\mathbf E(\mathbf r,\omega)
=\ii\omega\mu_0\int\dd^3r'\,
\mathbf G(\mathbf r,\mathbf r';\omega)
\mathbf j(\mathbf r',\omega).
\label{eq:feqo-green-source-field}
\end{equation}
若边界条件使齐次方程的解唯一，这就是
式~\eqref{eq:feqo-maxwell-operator} 的唯一场。
微分算子量纲为 \(\si{\per\square\metre}\)，
delta 分布为 \(\si{\per\cubic\metre}\)，
因此 \(\mathbf G\) 的量纲为 \(\si{\per\metre}\)。

\section{真空点源把归一化与辐射方向同时固定}
\label{sec:feqo-vacuum-dyadic}

在真空取 \(\mu_r=1,\boldsymbol\varepsilon=\mathbf I\)，
令 \(k=(\omega+\ii0)/c\) 和
\(R=|\mathbf r-\mathbf r'|\)。三维标量核
\(g_0(R)=\ee^{\ii kR}/(4\pi R)\) 在分布意义下满足
\((\nabla^2+k^2)g_0=-\delta^{(3)}(\mathbf r-\mathbf r')\)。
其电磁对应物为
\begin{equation}
\mathbf G_0(\mathbf r,\mathbf r';\omega)
=\left(\mathbf I+\frac{\nabla\nabla}{k^2}\right)
\frac{\ee^{\ii kR}}{4\pi R}.
\label{eq:feqo-vacuum-green}
\end{equation}
这里 \(\nabla\) 作用于第一个空间变量，\(\nabla\nabla g_0\)
必须作为分布理解，不能只在 \(R>0\) 处计算后漏掉接触项。

验证很短，但不能省略。记
\(\mathcal L_0=\nabla\times\nabla\times-k^2\mathbf I\)。
恒等式
\(\nabla\times\nabla\times(\mathbf I g_0)
=\nabla\nabla g_0-\mathbf I\nabla^2g_0\)
和 \(\nabla\times\nabla f=0\) 给
\[
\mathcal L_0\mathbf G_0
=-\mathbf I(\nabla^2+k^2)g_0
=\mathbf I\delta^{(3)}(\mathbf r-\mathbf r').
\]
所以式~\eqref{eq:feqo-vacuum-green} 连点源处的 delta 系数也正确。
\(\ee^{+\ii kR}\) 配合 \(\ee^{-\ii\omega t}\) 描述相位向外传播；
改成 \(\ee^{-\ii kR}\) 会改成向内来的解。

定义频域点电偶极矩 \(\mathbf d(\omega)\) 位于 \(\mathbf r_0\)，
其电流为
\(\mathbf j=-\ii\omega\mathbf d(\omega)
\delta^{(3)}(\mathbf r-\mathbf r_0)\)。
对应电荷密度是
\(\rho=-\mathbf d\cdot\nabla\delta^{(3)}(\mathbf r-\mathbf r_0)\)；
逐项代入 \(-\ii\omega\rho+\nabla\cdot\mathbf j=0\)
可见这个点源满足连续性方程，而不是凭空加入的电流。
代入式~\eqref{eq:feqo-green-source-field} 得
\[
\mathbf E(\mathbf r,\omega)
=\mu_0\omega^2
\mathbf G(\mathbf r,\mathbf r_0;\omega)\mathbf d(\omega).
\]
新定义立即算出一个实际场：给定偶极矩后，复杂结构的作用只需查
\(\mathbf G\) 的对应两点矩阵元。运动电子则把这个点源换成沿轨迹分布的电流，
因此会对 \(\mathbf r,\mathbf r'\) 作非局域采样。

真空偶极辐射给 Green 张量的系数一个独立检验。
在观测距离 \(R\gg\lambda=2\pi/k\) 处，
\(\nabla\nabla g_0=-k^2\hat{\mathbf R}\otimes
\hat{\mathbf R}\,g_0+O(R^{-2})\)，所以
\[
\mathbf G_0(\mathbf r,\mathbf r_0;\omega)
=\frac{\ee^{\ii kR}}{4\pi R}
(\mathbf I-\hat{\mathbf R}\otimes\hat{\mathbf R})
+O(R^{-2}).
\]
远场只保留横向偶极矩
\(\mathbf d_\perp=\mathbf d-
\hat{\mathbf R}(\hat{\mathbf R}\cdot\mathbf d)\)。
对于实场写法
\(\mathbf E_{\rm real}=\operatorname{Re}[
\boldsymbol{\mathcal E}\ee^{-\ii\omega t}]\)，
时间平均 Poynting 通量为
\(\frac12\varepsilon_0c|\boldsymbol{\mathcal E}|^2\)。
代入上面的远场，得到每单位立体角辐射功率
\[
\frac{\dd P}{\dd\Omega}
=\frac{\omega^4}{32\pi^2\varepsilon_0c^3}
|\mathbf d_\perp|^2.
\]
用 \(\int\hat R_i\hat R_j\,\dd\Omega
=4\pi\delta_{ij}/3\) 对全部方向积分，便有
\[
P_{\rm rad}
=\frac{\omega^4|\mathbf d|^2}{12\pi\varepsilon_0c^3}.
\]
这既检验了 \(4\pi R\) 的归一化，也说明“总损失”与某个有限
收集立体角内的“探测发光”不可能自动相等。

\section{模展开、因果性与损失通道}
\label{sec:feqo-green-spectral}

在无损、无色散的封闭腔中，取 \(\mu_r=1\) 和实正
\(\boldsymbol\varepsilon(\mathbf r)\)。
这里“横向”是相对于材料的 Gauss 条件
\(\nabla\cdot(\boldsymbol\varepsilon\mathbf f_\lambda)=0\)，
不必等于 \(\nabla\cdot\mathbf f_\lambda=0\)。
这些正常模满足
\[
\nabla\times\nabla\times\mathbf f_\lambda
=\frac{\omega_\lambda^2}{c^2}
\boldsymbol\varepsilon\mathbf f_\lambda,\qquad
\int\dd^3r\,
\mathbf f_\lambda^*\cdot\boldsymbol\varepsilon\mathbf f_{\lambda'}
=\delta_{\lambda\lambda'}.
\]
把 Green 方程在这些模式上投影，分母中的本征值差来自
左边的两项，得到横向部分
\begin{equation}
\mathbf G_{\rm T}(\mathbf r,\mathbf r';\omega)
=c^2\sum_\lambda
\frac{\mathbf f_\lambda(\mathbf r)
\otimes\mathbf f_\lambda^*(\mathbf r')}
{\omega_\lambda^2-(\omega+\ii0)^2}.
\label{eq:feqo-green-mode-expansion}
\end{equation}
这里 \(\mathbf u\otimes\mathbf v\) 表示把向量 \(\mathbf w\)
映为 \(\mathbf u(\mathbf v\cdot\mathbf w)\) 的线性算子，
因而上式确实是二阶张量，而不是两个向量的标量积。
这里并未声称该和式单独给出任意源的完整 Green 张量；
纵向静电响应须按 Gauss 定律另加。
例如真空 Fourier 空间的完整核还含
\(-c^2P_{\rm L}/(\omega+\ii0)^2\)，
其中 \(P_{\rm L}\) 把向量投影到波矢方向，
横向投影为 \(P_{\rm T}=I-P_{\rm L}\)。
对辐射的横向电流，式~\eqref{eq:feqo-green-mode-expansion}
与正常模描述逐项一致。

在 \(\omega>0\) 处，
\[
\operatorname{Im}
\frac1{\omega_\lambda^2-(\omega+\ii0)^2}
=\frac{\pi}{2\omega_\lambda}\delta(\omega-\omega_\lambda).
\]
所以 \(\operatorname{Im}\mathbf G_{\rm T}\) 在可激发的模式频率上
有正谱权重。更一般地，应把算符的虚部定义为
\(\operatorname{Im}\mathbf G=(\mathbf G-\mathbf G^\dagger)/(2\ii)\)；
在互易标量介质中它与常用的逐分量取虚部相容。

令 \(\langle\mathbf j,\mathbf G\mathbf j\rangle\)
表示对两个空间变量积分并对矢量指标缩并。
用式~\eqref{eq:feqo-green-source-field}，
外加源交给场的时间平均功率为
\[
P_{\rm src}
=-\frac12\operatorname{Re}\int\dd^3r\,
\mathbf j^*\cdot\mathbf E
=\frac{\omega\mu_0}{2}
\langle\mathbf j,\operatorname{Im}\mathbf G\,\mathbf j\rangle.
\]
被动材料要求 \(P_{\rm src}\ge0\)；这包括被材料吸收
和流向无穷远的辐射。若只取一组远场收集模式，
所得功率只是总量的一部分，不能把两者混为同一观测量。

被动性的条件还可以直接在 Maxwell 方程中检查，而不必仅凭
Green 张量的谱表示。在没有外加表面驱动的体积 \(V\) 内，
对频域方程使用恒等式
\(\nabla\cdot(\mathbf E\times\mathbf H^*)
=\mathbf H^*\cdot(\nabla\times\mathbf E)
-\mathbf E\cdot(\nabla\times\mathbf H^*)\)，
取实部并用散度定理，得到
\[
-\frac12\operatorname{Re}\int_V
\mathbf j^*\cdot\mathbf E\,\dd^3r
=\frac12\operatorname{Re}\int_{\partial V}
(\mathbf E\times\mathbf H^*)\cdot\mathbf n\,\dd A
+\frac{\omega\varepsilon_0}{2}\int_V
\mathbf E^*\cdot\operatorname{Im}_{\!H}
\boldsymbol\varepsilon\,\mathbf E\,\dd^3r
+\frac{\omega\mu_0}{2}\int_V
\operatorname{Im}\mu_r\,|\mathbf H|^2\,\dd^3r .
\]
这里 \(\operatorname{Im}_{\!H}\boldsymbol\varepsilon
=(\boldsymbol\varepsilon-\boldsymbol\varepsilon^\dagger)/(2\ii)\)；
磁导率在本式中取标量。右边依次是穿出边界的功率、
介电吸收和磁吸收。在 \(\omega>0\) 的约定下，
若材料局域且被动，后两项的系数非负；
把 \(V\) 扩到包含所有散射体且边界放在远场，
第一项也非负，于是前面的 Green 张量不等式才有物理依据。
有限体积内的穿出功率可带符号，不能单独把它叫作吸收。

互易性用的是另一种、\emph{不取复共轭}的积分。
若 \(\boldsymbol\varepsilon^\mathsf T
=\boldsymbol\varepsilon\)、\(\mu_r\) 是标量，
且两组解服从使边界积分相消的同一边界条件，
则对源 \(\mathbf j_1,\mathbf j_2\) 分别有
\(\mathcal L_\omega\mathbf E_a=\ii\omega\mu_0\mathbf j_a\)。
将第一式左乘 \(\mathbf E_2\)、第二式左乘
\(\mathbf E_1\) 后相减并积分，旋度项由分部积分
只留下已假定消失的边界项，材料项因转置对称相消，
故
\(\int\mathbf E_1\cdot\mathbf j_2\,\dd^3r
=\int\mathbf E_2\cdot\mathbf j_1\,\dd^3r\)。
依次把两源取成不同位置、不同方向的点电流，
并用式~\eqref{eq:feqo-green-source-field}，得到
\begin{equation}
G_{ij}(\mathbf r,\mathbf r';\omega)
=G_{ji}(\mathbf r',\mathbf r;\omega).
\label{eq:feqo-green-reciprocity}
\end{equation}
这里交换的是源点与观测点，同时转置矢量指标，
\emph{不是}对 Green 张量取复共轭。材料可以有吸收而仍满足
式~\eqref{eq:feqo-green-reciprocity}；非互易材料也仍可因果、
被动。因果、互易与被动因此是三个需要分别检验的条件。

\begin{exercise}
把 \(\mathbf j=-\ii\omega\mathbf d\,\delta_{\mathbf r_0}\)
代入功率公式，求它采样的 Green 张量矩阵元，
并解释为什么移动电子要采样两点而不是一点。
\end{exercise}


经典响应给出平均场和辐射功率。若光场最初有确定的光子数，
或者想知道交换后电子与光是否纠缠，只给平均场便不够。
电磁正常模的能量形式与谐振子相同，于是可以从熟悉的量子谐振子进入场量子化。

\section{一个电磁模式就是一个谐振子}
\label{sec:feqo-quantized-field}

先只考虑无损、无色散腔中的一个实驻波模式 \(\mathbf f(\mathbf r)\)，
频率为 \(\omega\)，并取上一节的归一化
\(\int\dd^3r\,\mathbf f\cdot\boldsymbol\varepsilon\mathbf f=1\)。
在没有自由电荷的辐射部分，选规范使标势为零，令
\(\mathbf A(\mathbf r,t)=Q(t)\mathbf f(\mathbf r)/\sqrt{\varepsilon_0}\)。
于是
\(\mathbf E=-\dot Q\,\mathbf f/\sqrt{\varepsilon_0}\)、
\(\mathbf B=Q\nabla\times\mathbf f/\sqrt{\varepsilon_0}\)。

把两场代入电磁能量
\(\frac12\int[\varepsilon_0\mathbf E\cdot
\boldsymbol\varepsilon\mathbf E+\mu_0^{-1}|\mathbf B|^2]\dd^3r\)。
电场项是 \(\dot Q^2/2\)。
对模式方程乘 \(\mathbf f\) 积分并分部积分；
腔边界条件使边界项为零，给
\(\int|\nabla\times\mathbf f|^2\dd^3r=\omega^2/c^2\)。
磁场项遂为 \(\omega^2Q^2/2\)。所以单模 哈密顿量 确实是
\[
H_\gamma=\frac12(P^2+\omega^2Q^2),\qquad P=\dot Q .
\]
这一计算决定了 \(Q\) 的单位和场的振幅，量子化时不再另猜模式体积因子。

\section{从正则对易关系到光子数态}

把 \(Q,P\) 换成算子并规定 \([Q,P]=\ii\hbar\)。
定义
\[
a=\sqrt{\frac{\omega}{2\hbar}}Q
+\frac{\ii}{\sqrt{2\hbar\omega}}P,\qquad
a^\dagger=\sqrt{\frac{\omega}{2\hbar}}Q
-\frac{\ii}{\sqrt{2\hbar\omega}}P.
\]
将两式相乘并用 \([Q,P]=\ii\hbar\)，得
\([a,a^\dagger]=1\)；反解为
\(Q=\sqrt{\hbar/(2\omega)}(a+a^\dagger)\)、
\(P=-\ii\sqrt{\hbar\omega/2}(a-a^\dagger)\)。
代回能量，
\begin{equation}
H_\gamma=\hbar\omega
\left(a^\dagger a+\frac12\right).
\label{eq:feqo-single-mode-hamiltonian}
\end{equation}
数算符 \(N=a^\dagger a\) 的本征态由
\(a|0\rangle=0\) 与
\(|n\rangle=(a^\dagger)^n|0\rangle/\sqrt{n!}\) 构造。
从 \([a,a^\dagger]=1\) 归纳得到
\(a|n\rangle=\sqrt n\,|n-1\rangle\)、
\(a^\dagger|n\rangle=\sqrt{n+1}\,|n+1\rangle\)。
固定光子数 \(n\) 时，吸收振幅含 \(\sqrt n\)，
发射振幅含 \(\sqrt{n+1}\)；后一个式子在 \(n=0\) 仍不为零。

Heisenberg 绘景中 \(a(t)=a\,\ee^{-\ii\omega t}\)。
由 \(\mathbf E=-\partial_t\mathbf A\) 得
\begin{equation}
\mathbf E^{(+)}(\mathbf r,t)
=\ii\sqrt{\frac{\hbar\omega}{2\varepsilon_0}}\,
\mathbf f(\mathbf r)a\,\ee^{-\ii\omega t},
\qquad
\mathbf E^{(-)}=(\mathbf E^{(+)})^\dagger,\quad
\mathbf E=\mathbf E^{(+)}+\mathbf E^{(-)}.
\label{eq:feqo-single-mode-field}
\end{equation}
这个系数的平方与 \(\hbar\omega\) 成正比，
正是每多一个量子需要增加的能量。
若模式有色散或损耗，本节使用的能量归一化不再直接适用；
那时应使用包含材料自由度的量子化或 Green 张量的连续环境表示。

开放空间没有离散的腔频率，却不需要重新猜一次量子场的系数。
先把真空装进边长 \(L\)、体积 \(V=L^3\) 的周期盒。
波矢 \(\mathbf k=2\pi\mathbf n/L\)，每个非零波矢有两个
单位横向偏振 \(\mathbf e_{\mathbf k\lambda}\)，满足
\(\mathbf k\cdot\mathbf e_{\mathbf k\lambda}=0\)。
离散算符满足
\([a_{\mathbf k\lambda},a_{\mathbf k'\lambda'}^\dagger]
=\delta_{\mathbf k\mathbf k'}\delta_{\lambda\lambda'}\)。
将平面波的 \(V^{-1/2}\) 归一化代入单模结果，
正频电场便是
\[
\mathbf E^{(+)}(\mathbf r,t)
=\ii\sum_{\mathbf k,\lambda}
\sqrt{\frac{\hbar c|\mathbf k|}{2\varepsilon_0 V}}\,
\mathbf e_{\mathbf k\lambda}a_{\mathbf k\lambda}
\ee^{\ii(\mathbf k\cdot\mathbf r-c|\mathbf k|t)}.
\]
这里对 \(\mathbf k\) 与 \(-\mathbf k\) 都求和，
厄米共轭再给负频部分；\(\mathbf k=0\) 的静电部分
不属于上述横向辐射模。
令 \(L\to\infty\) 时，每个动量格点占据
\((2\pi)^3/V\) 的波矢体积，因此
\(\sum_{\mathbf k}\to V(2\pi)^{-3}\int\dd^3k\)。
定义连续算符
\(a_\lambda(\mathbf k)
=\sqrt{V/(2\pi)^3}\,a_{\mathbf k\lambda}\)，
便有
\([a_\lambda(\mathbf k),a_{\lambda'}^\dagger(\mathbf k')]
=\delta_{\lambda\lambda'}\delta^{(3)}(\mathbf k-\mathbf k')\)，
而电场系数变为
\(\sqrt{\hbar c|\mathbf k|/[2\varepsilon_0(2\pi)^3]}\)。
这个极限说明连续谱中的 delta 归一化与盒中每模的
\(V^{-1/2}\) 是同一约定的两种写法。
单个确定 \(\mathbf k\) 的态同电子平面波一样不可归一化；
真正的一光子波包
\(\sum_\lambda\int F_\lambda(\mathbf k)
a_\lambda^\dagger(\mathbf k)|0\rangle\,\dd^3k\)
要求 \(\sum_\lambda\int|F_\lambda|^2\dd^3k=1\)。

这里也固定了实场与正频幅度的因子：
若一束经典单频场写成
\(\mathbf E_{\rm real}(t)
=\operatorname{Re}[\boldsymbol{\mathcal E}\ee^{-\ii\omega t}]\)，
则它的正频项是
\(\mathbf E^{(+)}=\boldsymbol{\mathcal E}
\ee^{-\ii\omega t}/2\)。
若把 \(\boldsymbol{\mathcal E}\) 直接当正频幅度，
后面沿轨迹的耦合会错一倍。

\section{相干态和热态由同一个数态基构造}
\label{sec:feqo-coherent-thermal}

\begin{definition}[相干态]
满足 \(a|\alpha\rangle=\alpha|\alpha\rangle\) 且范数为一的态，
称为振幅 \(\alpha\in\mathbb C\) 的相干态。
\end{definition}

写 \(|\alpha\rangle=\sum_{n\ge0}c_n|n\rangle\)。
比较 \(|n\rangle\) 的系数得到
\(\sqrt{n+1}c_{n+1}=\alpha c_n\)，
所以 \(c_n=c_0\alpha^n/\sqrt{n!}\)。
由 \(\sum_n|c_n|^2=|c_0|^2\ee^{|\alpha|^2}=1\)，
选择不影响物理的整体相位使 \(c_0>0\)，得到
\begin{equation}
|\alpha\rangle
=\ee^{-|\alpha|^2/2}
\sum_{n=0}^\infty\frac{\alpha^n}{\sqrt{n!}}|n\rangle,
\qquad
\Pr(N=n)=\ee^{-|\alpha|^2}\frac{|\alpha|^{2n}}{n!}.
\label{eq:feqo-coherent-state}
\end{equation}
这就是均值 \(\bar n=|\alpha|^2\) 的 Poisson 分布；
将求和指标平移一次、两次，分别得到
\(\langle N\rangle=|\alpha|^2\) 与
\(\langle N(N-1)\rangle=|\alpha|^4\)；
由 \(N^2=N(N-1)+N\) 得
\(\operatorname{Var}N=\bar n\)。
电场的平均正频部分为
\(\langle\mathbf E^{(+)}\rangle
=\ii\sqrt{\hbar\omega/(2\varepsilon_0)}
\mathbf f\alpha\ee^{-\ii\omega t}\)。
当 \(|\alpha|\) 增大、每光子耦合按 \(1/|\alpha|\) 缩放而
平均驱动保持固定时，光子数的相对涨落 \(1/|\alpha|\) 消失；
这才是后文所用的经典驱动极限。

若模式与温度 \(T\) 的热库平衡，Gibbs 权重按能量
\(E_n=\hbar\omega(n+\tfrac12)\) 分布。
令 \(x=\ee^{-\hbar\omega/(k_{\rm B}T)}\)，
配分函数中的零点能在归一化时相消，而
\(\sum_{n\ge0}x^n=(1-x)^{-1}\)。因此热态是
\begin{equation}
\rho_{\rm th}=(1-x)\sum_{n=0}^{\infty}
x^n|n\rangle\langle n|,
\qquad
\bar n=\frac{x}{1-x}
=\frac1{\ee^{\hbar\omega/(k_{\rm B}T)}-1}.
\label{eq:feqo-thermal-state}
\end{equation}
对 \(\sum_{n\ge0}x^n=(1-x)^{-1}\) 求一、二次导数，
得到 \(\sum n x^n=x/(1-x)^2\) 与
\(\sum n(n-1)x^n=2x^2/(1-x)^3\)。
代入 \(\langle N^2\rangle=\langle N(N-1)\rangle+\langle N\rangle\)，
便得数方差 \(\operatorname{Var}N=\bar n(1+\bar n)\)，
且 \(\langle a\rangle=0\)。
比较相干光和热光对电子的作用时，应固定相同的 \(\bar n\)
与电子耦合；否则输入平均能量不同也会改变边带宽度。

\section{电子电流怎样与光交换量子}
\label{sec:feqo-light-coupling}

上一章的 Dirac 哈密顿量 中，光场的矢势进入
\(-qc\boldsymbol\alpha\cdot\mathbf A\)。
本书先只讨论单电子空间，不预设量子场论的电子场算符。
对任意两条单电子 Dirac 波函数 \(\psi_i,\psi_f\)，定义
电荷流密度算符 \(\hat{\mathbf j}_{\rm ch}(\mathbf r)\)
的矩阵元为
\[
\langle\psi_f|\hat{\mathbf j}_{\rm ch}(\mathbf r)|\psi_i\rangle
=qc\,\psi_f^\dagger(\mathbf r)
\boldsymbol\alpha\psi_i(\mathbf r).
\]
特别地，它在同一态上的期望值正是上一章概率流的 \(q\) 倍。
在标势不参与辐射耦合的规范下，线性相互作用为
\begin{equation}
H_I(t)
=-\int\dd^3r\,
\hat{\mathbf j}_{\rm ch}(\mathbf r,t)
\cdot\hat{\mathbf A}(\mathbf r,t).
\label{eq:feqo-current-field-coupling}
\end{equation}
\(\hat{\mathbf j}_{\rm ch}\) 带电子电荷 \(q=-e\)，
单位为 \(\si{\ampere\per\square\metre}\)；
矢势单位为 \(\si{\volt\second\per\metre}\)，
故空间积分的单位确为焦耳。
不能与上一章的概率流混用。复杂材料中若另有标量势或纵向库模，
需把相应耦合一同保留。

为什么还需要纳米结构或近场？自由电子在真空中不能仅靠一次
发射变成另一自由电子加一个真实光子。
设初态 \((E_i,\mathbf p_i)\)，末电子
\((E_f,\mathbf p_f)\)，光子 \((\hbar\omega,\hbar\mathbf k)\)，
其中 \(\omega>0\)、\(|\mathbf k|=\omega/c\)。
若能量和动量都守恒，则
\(E_f=E_i-\hbar\omega\)、
\(\mathbf p_f=\mathbf p_i-\hbar\mathbf k\)。
把两式代入末电子的壳上条件
\(E_f^2-c^2|\mathbf p_f|^2=m^2c^4\)，
再减去初电子的同一条件，光子的
\(\hbar^2(\omega^2-c^2|\mathbf k|^2)\) 项恰好相消，留下
\[
E_i\omega=c^2\mathbf p_i\cdot\mathbf k.
\]
右边不超过 \(c|\mathbf p_i|\omega\)，
但有质量的电子满足 \(E_i>c|\mathbf p_i|\)，矛盾。
结构化近场可以携带不满足 \(|\mathbf k|=\omega/c\)
的纵向动量，并由材料整体接受剩余反冲；
这才使电子与光交换量子的通道打开。
这一点也可从轨迹相位看出。若场的一项沿电子方向写成
\(\ee^{\ii q_z z-\ii\omega t}\)，电子取 \(z=vt\) 时
采样的相位是 \(\ee^{\ii(q_zv-\omega)t}\)。
在低反冲的直线轨迹模型中，长距离相干累积需要
\(q_z\simeq\omega/v\)。
对有质量的电子 \(v<c\)，因此 \(q_z>\omega/c\)；
自由传播光子的纵向波数绝对值却不超过 \(\omega/c\)。
真空侧若仍满足局部 Maxwell 色散，横向波数便是
\(k_\perp=\ii\kappa\)、
\(\kappa=\sqrt{q_z^2-\omega^2/c^2}\)，
相应场随离开结构的距离按 \(\ee^{-\kappa x}\) 衰减。
这解释了为何相位匹配的分量通常是近场，而非远场平面波。

在相互作用绘景，有限时间 \(t_i\le t\le t_f\) 的演化是
\begin{equation}
U_I(t_f,t_i)=\mathcal T
\exp\!\left[-\frac{\ii}{\hbar}
\int_{t_i}^{t_f}H_I(t)\,\dd t\right].
\label{eq:feqo-interaction-propagator}
\end{equation}
时间排序 \(\mathcal T\) 把较晚时刻的算符放在左边；
它必需，因为不同时刻的 \(H_I(t)\) 未必对易。
弱耦合时展开到一次，
\[
U_I=I-\frac{\ii}{\hbar}\int_{t_i}^{t_f}H_I(t)\,\dd t
+O(H_I^2).
\]
若初态是 \(|i,n\rangle\)，一光子吸收的末态
\(|f,n-1\rangle\) 矩阵元含 \(\sqrt n\)，
发射到 \(|f,n+1\rangle\) 含 \(\sqrt{n+1}\)。
它们还各乘电子电流的空间重叠与有限时间的振荡积分；
下一章会从该积分算出谱线宽，而不会提前把有限时间换成能量 delta 函数。

\begin{exercise}
从 \(a|\alpha\rangle=\alpha|\alpha\rangle\) 逐项推出
式~\eqref{eq:feqo-coherent-state}，再求
\(\langle N(N-1)\rangle\) 以验证数方差。
\end{exercise}

\begin{exercise}
在固定平均光子数 \(\bar n\) 下，
分别计算数态 \(|n=\bar n\rangle\)（若 \(\bar n\) 为整数）、
相干态与热态的 \(\operatorname{Var}N\)。
说明“平均能量相同”并不推出“电子交换统计相同”。
\end{exercise}


环境的 Green 张量与量子正常模并不冲突：在无损封闭系统中，
二者通过谱展开逐项对应。材料存在吸收时，
仅保留那些无损腔模不足以描述吸收材料的自由度；
后文计算损失和噪声时将使用推迟响应及其连续谱。
